\section{Bibliografía}
% Formato APA 7. Todas las entradas tienen que estar citadas en el cuerpo del texto.
% La lista usa sangría francesa de 1.27 cm, como pide APA 7 (§2.12).
% Las cuatro entradas de herramientas (Anthropic, Harris et al., Hunter, Python Software
% Foundation) las cita Metodología. El resto se cita en el Marco teórico.
\begin{list}{}{%
  \setlength{\leftmargin}{1.27cm}%
  \setlength{\itemindent}{-1.27cm}%
  \setlength{\labelwidth}{0pt}%
  \setlength{\labelsep}{0pt}%
  \setlength{\itemsep}{0.6em}%
  \setlength{\parsep}{0pt}%
}
  \item[] Anthropic. (2026). \textit{Claude Code} (Fable 5.1) [Modelo de lenguaje grande]. Recuperado el 1 de noviembre de 2026, de \url{https://claude.ai/code}

  \item[] Burden, R. L., Faires, J. D., y Burden, A. M. (2016). \textit{Numerical analysis} (10.\textsuperscript{a} ed.). Cengage Learning.

  \item[] Cybersecurity and Infrastructure Security Agency. (s.f.). \textit{Reducing the significant risk of known exploited vulnerabilities}. Recuperado el 1 de noviembre de 2026, de \url{https://www.cisa.gov/known-exploited-vulnerabilities}

  \item[] Forum of Incident Response and Security Teams. (2024). \textit{Common Vulnerability Scoring System version 4.0: Specification document} (Revisión 1.2). \url{https://www.first.org/cvss/v4.0/specification-document}

  \item[] Harris, C. R., Millman, K. J., van der Walt, S. J., Gommers, R., Virtanen, P., Cournapeau, D., Wieser, E., Taylor, J., Berg, S., Smith, N. J., Kern, R., Picus, M., Hoyer, S., van Kerkwijk, M. H., Brett, M., Haldane, A., del Río, J. F., Wiebe, M., Peterson, P., \ldots{} Oliphant, T. E. (2020). Array programming with NumPy. \textit{Nature, 585}(7825), 357--362. \url{https://doi.org/10.1038/s41586-020-2649-2}

  \item[] Hunter, J. D. (2007). Matplotlib: A 2D graphics environment. \textit{Computing in Science \& Engineering, 9}(3), 90--95. \url{https://doi.org/10.1109/MCSE.2007.55}

  \item[] MITRE Corporation. (2021). \textit{CVE: Frequently asked questions} [Archivo del sitio anterior del programa CVE]. Recuperado el 1 de noviembre de 2026, de \url{https://www.cve.org/Resources/Media/Archives/OldWebsite/about/faqs.html}

  \item[] MITRE Corporation. (2026, 12 de mayo). \textit{Exploit public-facing application, technique T1190: Enterprise} (Versión 2.8). MITRE ATT\&CK. \url{https://attack.mitre.org/techniques/T1190/}

  \item[] National Institute of Standards and Technology. (s.f.). Vulnerability. En \textit{Computer Security Resource Center glossary}. Recuperado el 1 de noviembre de 2026, de \url{https://csrc.nist.gov/glossary/term/vulnerability}

  \item[] Python Software Foundation. (2024). \textit{The Python language reference} (Versión 3.12). Recuperado el 1 de noviembre de 2026, de \url{https://docs.python.org/3.12/reference/}

  \item[] Souppaya, M., y Scarfone, K. (2022). \textit{Guide to enterprise patch management planning: Preventive maintenance for technology} (NIST Special Publication 800-40, Rev. 4). National Institute of Standards and Technology. \url{https://doi.org/10.6028/NIST.SP.800-40r4}
\end{list}
